% =============================================================================
% Geo-SimCLR -- AMTH 5520 final-project Beamer deck
% =============================================================================
\documentclass[aspectratio=169,11pt]{beamer}

\usetheme{CambridgeUS}
\usecolortheme{beaver}
\setbeamertemplate{navigation symbols}{}     % cleaner footer (no tiny nav icons)
\setbeamertemplate{caption}[numbered]
\setbeamertemplate{itemize items}[circle]

\usepackage[T1]{fontenc}
\usepackage[utf8]{inputenc}
\usepackage{lmodern}
\usepackage{microtype}
\usepackage{graphicx}
\usepackage{booktabs}
\usepackage{amsmath}
\usepackage{amssymb}
\usepackage{xcolor}
\usepackage{hyperref}

% all figures live in charts/ relative to this .tex
\graphicspath{{charts/}}

% --- title-page metadata -----------------------------------------------------
\title[Geo-SimCLR]{Geo-SimCLR: Self-Supervised Contrastive Pre-training\\
                   of Sentinel-2 for U.S.\ Corn-Belt Yield Prediction}
\author[Ivan Petkov]{Ivan Petkov}
\institute[Yale]{%
  AMTH 5520 -- Final Project\\
  Yale University\\
  \vspace{2pt}
  \footnotesize Instructor: Prof.\ Smita Krishnaswamy
}
\date{May 2026}

% =============================================================================
\begin{document}
% =============================================================================

% ---------- Slide 1: title ---------------------------------------------------
\begin{frame}[plain]
  \titlepage
  \vfill
  \centering\scriptsize
  Code: \url{https://github.com/<TODO>/geo-simclr-cropnet}
  \quad$|$\quad
  Video: \url{https://youtu.be/<TODO>}
\end{frame}

% ---------- Slide 2: Motivation ----------------------------------------------
\section{Motivation}
\begin{frame}{Why this matters}
  \begin{itemize}\setlength\itemsep{6pt}
    \item \textbf{Crop-yield monitoring} drives agricultural planning, food
          security, monetary-policy inputs, and is now a recognised input to
          quantitative finance.
    \item \textbf{Sentinel-2 imagery is free \& global}: the bottleneck is
          \alert{representation}, not data acquisition.
    \item \textbf{Labels are scarce}: only $\sim$3\,k USDA county-year corn-yield
          labels across our 7-state Corn Belt sample, vs.\ $\sim$850\,k
          unlabelled tiles.
    \item \textbf{Recipe is broadly applicable}: the same gap exists in
          climate-change attribution, deforestation, biodiversity, urban-growth,
          and drought monitoring.
  \end{itemize}

  \vspace{6pt}
  \begin{block}{Goal}
    Learn a rich, transferable representation of multispectral satellite tiles
    \emph{from unlabelled imagery alone}, then probe it cheaply for yield.
  \end{block}
\end{frame}

% ---------- Slide 3: Background timeline -------------------------------------
\section{Background}
\begin{frame}{Three eras of crop-yield prediction from satellite imagery}
  \centering\small
  \renewcommand{\arraystretch}{1.35}
  \begin{tabular}{p{2.7cm} p{4.6cm} p{6.8cm}}
    \toprule
    \textbf{Era} & \textbf{Approach} & \textbf{Examples} \\
    \midrule
    Hand-crafted features &
      Spectral indices $\to$ tree / linear regressor &
      NDVI + LightGBM, Khaki et al.\ 2020 (CNN-RNN, tabular) \\
    Supervised CNN/RNN &
      End-to-end deep imagery models &
      ConvLSTM (Shi 2015), CNN-RNN, GNN-RNN (Fan 2022),
      MMST-ViT (Lin 2024, multi-modal) \\
    Self-supervised &
      Contrastive / MAE pre-training, then probe &
      SimCLR (Chen 2020), DINOv2, MAE / SatMAE,
      Prithvi-EO 2.0, SeCo, GASSL\\
    \midrule
    \alert{Our work} &
      \alert{SimCLR with temporal positives} &
      \alert{Geo-SimCLR (this project)} \\
    \bottomrule
  \end{tabular}
\end{frame}

% ---------- Slide 4: Data overview -------------------------------------------
\section{Data}
\begin{frame}{Data: 7 Corn-Belt states, 6 years, biweekly Sentinel-2}
  \centering
  \includegraphics[width=0.95\linewidth,height=0.78\textheight,keepaspectratio]%
                  {fig03_data_summary.png}

  \vspace{4pt}
  \tiny CropNet release of Sentinel-2 imagery -- agriculture false-colour
  bands B02 / B08 / B11 -- biweekly Apr--Sep, 224$\times$224 px @ 40 m/pixel
  ($9 \times 9$\,km tiles).
\end{frame}

% ---------- Slide 5: Phenology + yield variance ------------------------------
\begin{frame}{Phenology window and yield variance}
  \begin{columns}[T]
    \begin{column}{0.55\textwidth}
      \centering
      \includegraphics[width=\linewidth,height=0.55\textheight,keepaspectratio]%
                      {figA2_phen_stage.png}
      \\[2pt]\tiny USDA Handbook 628 reference phenology -- 7-state coverage
      of the full April--September growing window.
    \end{column}
    \begin{column}{0.43\textwidth}
      \centering
      \includegraphics[width=\linewidth,height=0.55\textheight,keepaspectratio]%
                      {figA3_yield_distribution.png}
      \\[2pt]\tiny Per-state yield boxplots (early 3-state snapshot;
      full 7-state numbers in the report).
    \end{column}
  \end{columns}
\end{frame}

% ---------- Slide 6: Pipeline schematic --------------------------------------
\section{Method}
\begin{frame}{Geo-SimCLR pipeline}
  \centering
  \includegraphics[width=0.95\linewidth,height=0.80\textheight,keepaspectratio]%
                  {fig01_pipeline.png}
\end{frame}

% ---------- Slide 7: Augmentation examples on real tiles ---------------------
\begin{frame}{What the contrastive task looks like on real tiles}
  \centering
  \includegraphics[width=0.95\linewidth,height=0.83\textheight,keepaspectratio]%
                  {fig02_augmentation_examples.png}
\end{frame}

% ---------- Slide 8: Method details (loss + hyperparameters) -----------------
\begin{frame}{NT-Xent loss and key training settings}
  \vspace{-4pt}
  \begin{block}{NT-Xent loss (Chen et al.\ 2020)}
    \vspace{-8pt}
    \[
      \mathcal{L}_{i \to j} \;=\; -\log
        \frac{\exp(z_i^\top z_j / \tau)}
             {\sum_{k=1}^{2N} \mathbf{1}_{[k\neq i]} \exp(z_i^\top z_k / \tau)}
    \]
    \vspace{-10pt}
    \centering\scriptsize
    $\ell_2$-normalised projections $z_i$, in-batch negatives,
    temperature $\tau = 0.20$.
  \end{block}

  \vspace{-2pt}
  \begin{columns}[T]
    \begin{column}{0.50\textwidth}
      \footnotesize\textbf{Architecture}
      \begin{itemize}\footnotesize\setlength\itemsep{0pt}
        \item ResNet-18 from scratch, 11.2\,M params
        \item Projector MLP $512 \to 256 \to 128$ (dropped after SSL)
      \end{itemize}
      \footnotesize\textbf{Optimisation}
      \begin{itemize}\footnotesize\setlength\itemsep{0pt}
        \item AdamW, lr $10^{-3}$, wd $10^{-4}$
        \item 5-ep warm-up $\to$ cosine, 50 epochs
        \item Batch 1024, AMP, $\sim$5\,h on 2 GPUs
      \end{itemize}
    \end{column}
    \begin{column}{0.50\textwidth}
      \footnotesize\textbf{Pretext-task design}
      \begin{itemize}\footnotesize\setlength\itemsep{0pt}
        \item Positive = same field, biweek $t \pm 1$ ($\sim$14 d)
        \item CRD-balanced anchors (61 unique CRDs)
        \item $\sim$645\,k anchor pairs / epoch
      \end{itemize}
      \footnotesize\textbf{Augmentations (per view, indep.)}
      \begin{itemize}\footnotesize\setlength\itemsep{0pt}
        \item h/v flip, rot90, crop 160--208 $\to$ 224
        \item sharpness, intensity jitter, Gaussian noise
        \item \alert{No} spectral dropout, RGB jitter, mixup
      \end{itemize}
    \end{column}
  \end{columns}
\end{frame}

% ---------- Slide 9: Evaluation methodology ----------------------------------
\begin{frame}{Evaluation methodology}
  \begin{columns}[T]
    \begin{column}{0.50\textwidth}
      \textbf{Splits}
      \begin{itemize}\small\setlength\itemsep{3pt}
        \item \textbf{Time split:} train 2017--2020 / val 2021 /
              test 2022. Strict no-look-ahead in-season-forecast benchmark.
        \item \textbf{County split:} random 20\,\% of FIPS held out across
              all years. Spatial generalisation test.
      \end{itemize}
    \end{column}
    \begin{column}{0.50\textwidth}
      \textbf{Probes on frozen features}
      \begin{itemize}\small\setlength\itemsep{3pt}
        \item Mean-pool tile features per (\textit{fips, year}) into one
              512-d vector.
        \item Probe: \textbf{Ridge} ($\alpha\!=\!1$) or \textbf{LightGBM}
              (300 trees, early stopping on 2021).
        \item Optional state one-hot + year scalar: only on county split
              (year extrapolates badly out-of-sample).
        \item Bootstrap 95\,\% CIs from $n\!=\!200$ paired resamples.
      \end{itemize}
    \end{column}
  \end{columns}

  \vspace{4pt}
  \begin{block}{Why split-aware probes}
    For \emph{every} encoder, Ridge wins time-split and LightGBM with state+year
    wins county-split. Reported headline cells are the empirical winner per
    encoder (full grid in the report's Appendix).
  \end{block}
\end{frame}

% ---------- Slide 10: Baselines ----------------------------------------------
\begin{frame}{Baselines: imagery-only, fixed input modality}
  \begin{itemize}\setlength\itemsep{6pt}
    \item \textbf{Hand-crafted NDVI} (6 specifications) $\to$ Ridge / LightGBM.
          Strongest: state + year + biweekly NDVI series (LightGBM, 20 features).
    \item \textbf{Random ResNet-18}: same architecture as Geo-SimCLR, never trained.
          Sanity baseline -- mean-pooled random conv features carry colour /
          texture statistics correlated with vegetation.
    \item \textbf{DINOv2-base} (frozen, 85.7\,M params, 768-d).
          Trained on 142\,M natural RGB images; AG bands fed as RGB slots.
    \item \textbf{Prithvi-EO 2.0} (frozen, 303.9\,M params, 1024-d).
          EO-specialised but expects 6-band HLS; we channel-pad the 3 missing.
    \item \textbf{ConvLSTM} end-to-end (12.4\,M params).
          Per-frame ResNet-18 + LSTM(512$\to$256) + MLP head, $z$-scored
          target, image augmentation across all 12 biweek frames.
  \end{itemize}

  \vspace{2pt}
  \scriptsize\itshape We deliberately do not chase MMST-ViT: it is multi-modal
  and the CropNet ablation attributes most of its gain to the WRF-HRRR weather
  modality, not the imagery.
\end{frame}

% ---------- Slide 11: Headline results ---------------------------------------
\section{Results}
\begin{frame}{Headline result: best single encoder on both splits}
  \centering
  \includegraphics[width=0.92\linewidth,height=0.62\textheight,keepaspectratio]%
                  {fig05_encoder_comparison.png}

  \vspace{2pt}
  \begin{block}{Headline numbers}
    \centering\small Geo-SimCLR test
    $R^2 = \textbf{0.669}$ (time, Ridge) and $\textbf{0.660}$ (county, LightGBM);
    +$0.118\,R^2$ over supervised end-to-end ConvLSTM on time.
  \end{block}
\end{frame}

% ---------- Slide 12: Why it works -------------------------------------------
\begin{frame}{Reading the ranking}
  \begin{itemize}\setlength\itemsep{8pt}
    \item \textbf{Geo-SimCLR $\succ$ ConvLSTM end-to-end} ($+0.118\,R^2$
          time) -- SSL pretext sees $\sim$6.4\,M (anchor, positive,
          1023-negs) tuples $\times$ 50 epochs over 850\,k unlabelled tiles
          vs.\ ConvLSTM's $\sim$2{,}100 supervised yield labels. Pure
          \alert{data-efficiency} win.
    \item \textbf{Geo-SimCLR $\succ$ Prithvi-EO 2.0} (despite 27$\times$ more
          parameters) -- Prithvi expects 6-band HLS; we channel-pad the
          missing 3. Effective rank on padded input collapses to
          $\sim$$1.7 / 1024$ -- documented \alert{baseline limitation}, not a
          claim about Prithvi-the-model.
    \item \textbf{Geo-SimCLR $\approx$ DINOv2-base} (within $\sim$$0.06\,R^2$
          at $7.6\times$ fewer parameters) -- large general-purpose SSL
          transfers surprisingly well to satellite RGB-like inputs.
          Geo-SimCLR's edge comes from the temporal-positive pretext task
          being domain-aware.
  \end{itemize}
\end{frame}

% ---------- Slide 13: Training trajectory ------------------------------------
\begin{frame}{SSL training trajectory (and the first warning sign)}
  \centering
  \includegraphics[width=0.95\linewidth,height=0.78\textheight,keepaspectratio]%
                  {fig06_training_curves.png}

  \vspace{2pt}
  \scriptsize Loss converges cleanly; encoder norm is stable (no zero-collapse).
  But mean positive-pair cosine similarity climbs to $\sim$$0.95$ -- adjacent-biweek
  positives are intrinsically very similar.
\end{frame}

% ---------- Slide 14: M-PHATE ------------------------------------------------
\begin{frame}{What did the encoder learn? -- M-PHATE}
  \centering
  \includegraphics[width=0.95\linewidth,height=0.78\textheight,keepaspectratio]%
                  {fig07_mphate.png}

  \vspace{2pt}
  \scriptsize \textbf{State coloring (a):} clean state-coherent clusters --
  geography is encoded.\quad
  \textbf{Biweek coloring (b):} smooth but tightly overlapping --
  seasonal phenology only weakly encoded.
\end{frame}

% ---------- Slide 15: Honest collapse diagnostic -----------------------------
\begin{frame}{Honest collapse diagnostic}
  \centering
  \includegraphics[width=0.95\linewidth,height=0.62\textheight,keepaspectratio]%
                  {fig08_collapse.png}

  \vspace{2pt}
  \begin{block}{Three numbers that quantify it}
    \centering\small
    Effective rank \textbf{82.3\,/\,512} \,$|$\,
    mean post-$\ell_2$ cosine sim \textbf{+0.228} \,$|$\,
    each county's 12 biweek points span only \textbf{4--6\,\%} of the global
    pairwise distance.
  \end{block}
\end{frame}

% ---------- Slide 16: GradCAM ------------------------------------------------
\begin{frame}{Pixel attribution via GradCAM}
  \centering
  \includegraphics[width=0.95\linewidth,height=0.74\textheight,keepaspectratio]%
                  {fig09_gradcam.png}

  \vspace{2pt}
  \scriptsize Read in light of the collapse: encoder finds \emph{where} the
  cropland is -- consistent across high- and low-yield extremes -- not
  \emph{how that crop is currently doing}.
\end{frame}

% ---------- Slide 17: Limitations + future work ------------------------------
\section{Discussion}
\begin{frame}{Limitations and future work}
  \begin{columns}[T]
    \begin{column}{0.48\textwidth}
      \textbf{Limitations}
      \begin{itemize}\small\setlength\itemsep{3pt}
        \item Partial dimensional collapse along county-identity directions
              ($R_{\mathrm{eff}}\!=\!82/512$).
        \item Phenology only weakly encoded -- our headline claim is
              \emph{strong geography}, not \emph{strong phenology}.
        \item Small backbone (ResNet-18, 11\,M params) by SSL-paper standards.
        \item Evaluation restricted to 7 states / 6 years / corn only.
      \end{itemize}
    \end{column}
    \begin{column}{0.52\textwidth}
      \textbf{Future work (priority order)}
      \begin{enumerate}\small\setlength\itemsep{3pt}
        \item \alert{Harder positives + stronger augmentations}
              ($\pm 3$--$5$ biweeks; tighter crop ranges).
        \item \alert{Decorrelation regulariser} (Barlow-Twins / VICReg term
              added to NT-Xent).
        \item Head-to-head against SeCo, GASSL, GeoSSL, SatMAE under the same
              probe protocol.
        \item 4-channel AG+NDVI input;\; partial fine-tuning;\;
              fine-tune on USDA Weekly Crop Progress.
        \item Larger backbone, more crops (soybean transfer), weather ablation,
              cross-country generalisation.
      \end{enumerate}
    \end{column}
  \end{columns}
\end{frame}

% ---------- Slide 18: Thank-you / references ---------------------------------
\begin{frame}[plain]{Thank you  \quad{\normalfont\small Questions?}}
  \vspace{6pt}
  \footnotesize
  \textbf{Full report:} accompanying NeurIPS-style paper (this folder).\\
  \textbf{Code:} \url{https://github.com/<TODO>/geo-simclr-cropnet}\\
  \textbf{Video:} \url{https://youtu.be/<TODO>}

  \vspace{14pt}
  \textbf{Selected references}
  \scriptsize
  \begin{itemize}\setlength\itemsep{1pt}
    \item Chen et al.\ \textit{A Simple Framework for Contrastive Learning of
          Visual Representations.} ICML 2020. (SimCLR)
    \item Lin et al.\ \textit{An Open and Large-Scale Dataset for Multi-Modal
          Climate Change-Aware Crop Yield Predictions.} KDD 2024. (CropNet)
    \item Oquab et al.\ \textit{DINOv2: Learning Robust Visual Features
          without Supervision.} TMLR 2024.
    \item Jakubik et al.\ \textit{Prithvi-EO-2.0: A Versatile Multi-Temporal
          Foundation Model for Earth Observation.} IBM-NASA 2024.
    \item Shi et al.\ \textit{Convolutional LSTM Network: A Machine Learning
          Approach for Precipitation Nowcasting.} NeurIPS 2015.
    \item Gigante et al.\ \textit{Visualizing the PHATE of Neural Networks.}
          NeurIPS 2019. (M-PHATE)
    \item Jing et al.\ \textit{Understanding Dimensional Collapse in
          Contrastive Self-Supervised Learning.} ICLR 2022.
    \item Selvaraju et al.\ \textit{Grad-CAM: Visual Explanations from Deep
          Networks via Gradient-Based Localization.} ICCV 2017.
  \end{itemize}
\end{frame}

\end{document}
